\documentclass[10pt]{extarticle}
\usepackage[a4paper,margin=1.5cm]{geometry}
\usepackage[T1]{fontenc}
\usepackage{amsmath,amssymb}
\usepackage{graphicx}
\usepackage[export]{adjustbox}
\usepackage{multicol}
\usepackage{enumitem}
\usepackage{booktabs,array,tabularx}
\setlength{\columnsep}{0.7cm}
\setlength{\parindent}{0pt}
\setlength{\parskip}{3pt}
\pagestyle{empty}
\makeatletter
% figure registry (figdefs.tex is generated by the extract step)
\newcommand{\qfigdef}[4]{\@namedef{qf@file@#1}{#2}\@namedef{qf@w@#1}{#3}\@namedef{qf@wide@#1}{#4}}
\newcommand{\qf@check}[1]{\@ifundefined{qf@file@#1}{\PackageError{score2}{Unknown figure id #1}{}}{}}
\newcommand{\qf@img}[4][]{\includegraphics[width=#2\dimexpr\@nameuse{qf@w@#4}pt\relax,max width=#3,max height=0.45\textheight,#1]{\@nameuse{qf@file@#4}}}
% Full-width questions: finalize wraps every question that contains a wide figure in \qfullbegin...\qfullend,
% so the WHOLE question (stem, figures, options) is set across the page in reading order, then the two columns
% resume. The question's head - its number, stem and first figure (\qfullsplit marks the end of the head) -
% is kept on one page, so a number is never left at a page foot with its figure overleaf; the rest may break
% across pages normally, so no big gaps are left. Without a \qfullsplit, the whole question is kept together
% if it is under 45% of a page. Never write these yourself.
\newif\ifqf@full
\newif\ifqf@boxed
\newsavebox{\qf@box}
\newcommand{\qf@open}{\global\qf@fulltrue\par\global\qf@boxedtrue
  \setbox\qf@box\vbox\bgroup\hsize\textwidth\linewidth\textwidth}
\newcommand{\qf@emit}[1]{\par\egroup\global\qf@boxedfalse
  \ifdim\dimexpr\ht\qf@box+\dp\qf@box\relax<#1\textheight\noindent\usebox{\qf@box}\par\else\unvbox\qf@box\fi}
\newcommand{\qfullsplit}{\ifqf@boxed\qf@emit{0.85}\fi}
\newcommand{\qf@close}{\ifqf@boxed\qf@emit{0.45}\fi\global\qf@fullfalse}
\newcommand{\qfullbegin}{\end{multicols}\qf@open}
\newcommand{\qfullend}{\qf@close\begin{multicols}{2}}
\newcommand{\qfullstart}{\qf@open}
\newcommand{\qfulldone}{\qf@close}
% \qfig[scale]{id}: block figure, centred, directly where it is written (size set by the tool).
% Fallback if a wide figure is not inside a full-width question: break out of the columns in place - no floats.
\newcommand{\qfig}[2][1]{\qf@check{#2}%
  \ifqf@full\qf@blk{#1}{#2}\else
  \ifnum\@nameuse{qf@wide@#2}=1 \expandafter\@firstoftwo\else\expandafter\@secondoftwo\fi
  {\end{multicols}\begin{center}\qf@img{#1}{\textwidth}{#2}\end{center}\begin{multicols}{2}}%
  {\qf@blk{#1}{#2}}\fi}
\newcommand{\qf@blk}[2]{\par\smallskip{\centering\qf@img{#1}{\linewidth}{#2}\par}\smallskip}
% \qfigi[scale]{id}: inline figure (inside option lists / tables / minipages); never breaks columns.
\newcommand{\qfigi}[2][1]{\qf@check{#2}\qf@img[valign=c]{#1}{\linewidth}{#2}}
\makeatother
% structure macros - lint reads these to check completeness and figure placement
\newcommand{\papersec}[1]{\par\bigskip{\centering\Large\bfseries PAPER-#1\par}\medskip}
\newcommand{\partsec}[2]{\par\medskip\noindent{\large\bfseries PART-#1: #2}\par\smallskip}
\newcommand{\sect}[1]{\par\smallskip\noindent\textbf{#1}\par}
\newcommand{\qn}[1]{\par\smallskip\noindent\textbf{#1.}~}
\newcommand{\paraq}[2]{\par\smallskip\noindent\textbf{Paragraph for Questions #1 and #2}\par}
\newenvironment{qoptions}{\begin{itemize}[leftmargin=1.8em,itemsep=1pt,topsep=2pt,label={}]}{\end{itemize}}
\newcommand{\opt}[2]{\item[(#1)] #2}
% figure registry (generated by score2_tools.py finalize)
\qfigdef{p1_phy_s1_q01_1}{Score Advanced - Test 33 (12 May 2024)-img/p1_phy_s1_q01_1.png}{83.5}{0}
\qfigdef{p1_phy_s1_q01_3}{Score Advanced - Test 33 (12 May 2024)-img/p1_phy_s1_q01_3.png}{125.7}{0}
\qfigdef{p1_phy_s1_q04_1}{Score Advanced - Test 33 (12 May 2024)-img/p1_phy_s1_q04_1.png}{117.2}{0}
\qfigdef{p1_phy_s2_q07_1}{Score Advanced - Test 33 (12 May 2024)-img/p1_phy_s2_q07_1.png}{101.8}{0}
\qfigdef{p1_phy_s3_q08_1}{Score Advanced - Test 33 (12 May 2024)-img/p1_phy_s3_q08_1.png}{76.0}{0}
\qfigdef{p1_phy_s3_q10_1}{Score Advanced - Test 33 (12 May 2024)-img/p1_phy_s3_q10_1.png}{86.2}{0}
\qfigdef{p1_phy_s3_q11_P}{Score Advanced - Test 33 (12 May 2024)-img/p1_phy_s3_q11_P.png}{70.3}{0}
\qfigdef{p1_phy_s3_q11_Q}{Score Advanced - Test 33 (12 May 2024)-img/p1_phy_s3_q11_Q.png}{91.8}{0}
\qfigdef{p1_phy_s3_q11_R}{Score Advanced - Test 33 (12 May 2024)-img/p1_phy_s3_q11_R.png}{66.7}{0}
\qfigdef{p1_phy_s3_q11_S}{Score Advanced - Test 33 (12 May 2024)-img/p1_phy_s3_q11_S.png}{96.8}{0}
\qfigdef{p1_phy_s4_q04_1}{Score Advanced - Test 33 (12 May 2024)-img/p1_phy_s4_q04_1.png}{86.2}{0}
\qfigdef{p1_phy_s4_q06_1}{Score Advanced - Test 33 (12 May 2024)-img/p1_phy_s4_q06_1.png}{69.3}{0}
\qfigdef{p1_chem_s1_q03_1}{Score Advanced - Test 33 (12 May 2024)-img/p1_chem_s1_q03_1.png}{165.7}{0}
\qfigdef{p1_chem_s1_q04_1}{Score Advanced - Test 33 (12 May 2024)-img/p1_chem_s1_q04_1.png}{262.0}{1}
\qfigdef{p1_chem_s1_q04_O}{Score Advanced - Test 33 (12 May 2024)-img/p1_chem_s1_q04_O.png}{308.8}{1}
\qfigdef{p1_chem_s2_q05_1}{Score Advanced - Test 33 (12 May 2024)-img/p1_chem_s2_q05_1.png}{102.2}{0}
\qfigdef{p1_chem_s2_q06_1}{Score Advanced - Test 33 (12 May 2024)-img/p1_chem_s2_q06_1.png}{73.0}{0}
\qfigdef{p1_chem_s3_q08_1}{Score Advanced - Test 33 (12 May 2024)-img/p1_chem_s3_q08_1.png}{131.7}{0}
\qfigdef{p1_chem_s3_q10_P}{Score Advanced - Test 33 (12 May 2024)-img/p1_chem_s3_q10_P.png}{107.0}{0}
\qfigdef{p1_chem_s3_q10_Q}{Score Advanced - Test 33 (12 May 2024)-img/p1_chem_s3_q10_Q.png}{93.2}{0}
\qfigdef{p1_chem_s3_q10_R}{Score Advanced - Test 33 (12 May 2024)-img/p1_chem_s3_q10_R.png}{76.3}{0}
\qfigdef{p1_chem_s3_q10_S}{Score Advanced - Test 33 (12 May 2024)-img/p1_chem_s3_q10_S.png}{79.5}{0}
\qfigdef{p1_chem_s4_q05_1}{Score Advanced - Test 33 (12 May 2024)-img/p1_chem_s4_q05_1.png}{241.3}{0}
\qfigdef{p2_phy_s1_q01_1}{Score Advanced - Test 33 (12 May 2024)-img/p2_phy_s1_q01_1.png}{94.2}{0}
\qfigdef{p2_phy_s1_q03_1}{Score Advanced - Test 33 (12 May 2024)-img/p2_phy_s1_q03_1.png}{106.2}{0}
\qfigdef{p2_phy_s1_q04_1}{Score Advanced - Test 33 (12 May 2024)-img/p2_phy_s1_q04_1.png}{55.0}{0}
\qfigdef{p2_phy_s1_q04_A}{Score Advanced - Test 33 (12 May 2024)-img/p2_phy_s1_q04_A.png}{59.3}{0}
\qfigdef{p2_phy_s1_q04_B}{Score Advanced - Test 33 (12 May 2024)-img/p2_phy_s1_q04_B.png}{59.3}{0}
\qfigdef{p2_phy_s1_q04_C}{Score Advanced - Test 33 (12 May 2024)-img/p2_phy_s1_q04_C.png}{59.3}{0}
\qfigdef{p2_phy_s1_q04_D}{Score Advanced - Test 33 (12 May 2024)-img/p2_phy_s1_q04_D.png}{59.3}{0}
\qfigdef{p2_phy_s2_q06_1}{Score Advanced - Test 33 (12 May 2024)-img/p2_phy_s2_q06_1.png}{138.3}{0}
\qfigdef{p2_phy_s2_q08_1}{Score Advanced - Test 33 (12 May 2024)-img/p2_phy_s2_q08_1.png}{85.8}{0}
\qfigdef{p2_phy_s2_q10_1}{Score Advanced - Test 33 (12 May 2024)-img/p2_phy_s2_q10_1.png}{37.4}{0}
\qfigdef{p2_phy_s3_q06_1}{Score Advanced - Test 33 (12 May 2024)-img/p2_phy_s3_q06_1.png}{65.3}{0}
\qfigdef{p2_phy_s3_q08_1}{Score Advanced - Test 33 (12 May 2024)-img/p2_phy_s3_q08_1.png}{91.5}{0}
\qfigdef{p2_chem_s1_q01_1}{Score Advanced - Test 33 (12 May 2024)-img/p2_chem_s1_q01_1.png}{155.7}{0}
\qfigdef{p2_chem_s1_q04_1}{Score Advanced - Test 33 (12 May 2024)-img/p2_chem_s1_q04_1.png}{245.2}{0}
\qfigdef{p2_chem_s2_q05_S}{Score Advanced - Test 33 (12 May 2024)-img/p2_chem_s2_q05_S.png}{139.5}{0}
\qfigdef{p2_chem_s2_q05_P}{Score Advanced - Test 33 (12 May 2024)-img/p2_chem_s2_q05_P.png}{38.0}{0}
\qfigdef{p2_chem_s2_q05_Y}{Score Advanced - Test 33 (12 May 2024)-img/p2_chem_s2_q05_Y.png}{20.3}{0}
\qfigdef{p2_chem_s2_q05_X}{Score Advanced - Test 33 (12 May 2024)-img/p2_chem_s2_q05_X.png}{19.3}{0}
\qfigdef{p2_chem_s2_q06_a}{Score Advanced - Test 33 (12 May 2024)-img/p2_chem_s2_q06_a.png}{55.0}{0}
\qfigdef{p2_chem_s2_q06_b}{Score Advanced - Test 33 (12 May 2024)-img/p2_chem_s2_q06_b.png}{67.6}{0}
\qfigdef{p2_chem_s2_q06_c}{Score Advanced - Test 33 (12 May 2024)-img/p2_chem_s2_q06_c.png}{36.7}{0}
\qfigdef{p2_chem_s2_q06_d}{Score Advanced - Test 33 (12 May 2024)-img/p2_chem_s2_q06_d.png}{53.0}{0}
\qfigdef{p2_chem_s3_q06_A}{Score Advanced - Test 33 (12 May 2024)-img/p2_chem_s3_q06_A.png}{119.7}{0}
\qfigdef{p2_chem_s3_q06_8}{Score Advanced - Test 33 (12 May 2024)-img/p2_chem_s3_q06_8.png}{226.2}{0}
\qfigdef{p2_chem_s3_q08_1}{Score Advanced - Test 33 (12 May 2024)-img/p2_chem_s3_q08_1.png}{36.3}{0}
\raggedcolumns
\makeatletter\@beginparpenalty=10000\makeatother
\newcolumntype{L}[1]{>{\raggedright\arraybackslash}p{#1}}
\begin{document}
\begin{center}{\Large\bfseries Score Advanced --- Test 33 (12 May 2024)}\end{center}
\begin{multicols}{2}
\raggedcolumns
\papersec{1}
\partsec{1}{PHYSICS}
\sect{SECTION-I(i)}
\qn{1} A block is resting on a horizontal plate in the XY plane and co-efficient of friction between block and the plate is $\mu$. The plate begins to move with velocity v = bt$^2$ in X-direction. At what time will the block starts sliding from plate ?\par
\qfig{p1_phy_s1_q01_1}
\begin{qoptions}
\opt{A}{$\dfrac{\mu b}{g}$}
\opt{B}{$\dfrac{\mu bg}{2}$}
\opt{C}{$\dfrac{\mu g}{b}$}
\opt{D}{$\dfrac{\mu g}{2b}$}
\end{qoptions}
\qn{2} Monochromatic light of wavelength $\lambda$ emerging from slit S illuminates slits S$_1$ and S$_2$ which are placed with respect to S as shown in figure. The distances x and D are large compared to the separation d between the slits. If x = D/2, the minimum value of d so that there is a dark fringe at the centre P of the screen is :-\par
\qfig{p1_phy_s1_q01_3}
\begin{qoptions}
\opt{A}{$\sqrt{\dfrac{\lambda D}{3}}$}
\opt{B}{$\sqrt{\dfrac{2\lambda D}{3}}$}
\opt{C}{$\sqrt{\lambda D}$}
\opt{D}{$2\sqrt{\dfrac{\lambda D}{3}}$}
\end{qoptions}
\qn{3} An electromagnetic wave of frequency f = 3.0 MHz passes from vacuum into a dielectric medium with permittivity e = 4.0. Then :-
\begin{qoptions}
\opt{A}{wavelength is doubled and the frequency remains unchanged}
\opt{B}{wavelength is doubled and frequency becomes half.}
\opt{C}{wavelength is halved and frequency remains unchanged.}
\opt{D}{wavelength and frequency both remain unchanged.}
\end{qoptions}
\qn{4} Switch is in position A for a long time. It is suddenly switched to the side B. Mass of rod is m and length is $\ell$. Maximum speed of rod will be :-\par
\qfig{p1_phy_s1_q04_1}
\begin{qoptions}
\opt{A}{$\left(\dfrac ER\right)\sqrt{\dfrac Lm}$}
\opt{B}{$\left(\dfrac ER\right)\sqrt{\dfrac mL}$}
\opt{C}{$2\left(\dfrac ER\right)\sqrt{\dfrac Lm}$}
\opt{D}{$2\left(\dfrac ER\right)\sqrt{\dfrac mL}$}
\end{qoptions}
\sect{SECTION-I(ii)}
\qn{5} The current density in a wire of radius a varies with radial distance r as J = kr$^2$, where k is a constant.
\begin{qoptions}
\opt{A}{Total current passing through the cross section of the wire is I = $\dfrac{\pi ka^4}2$.}
\opt{B}{Total current passing through the cross section of the wire is I = $\dfrac{3\pi ka^3}2$}
\opt{C}{The field at a distance r $>$ a is B = $\dfrac{\mu_0ka^4}{4r}$}
\opt{D}{The field at a distance r $<$ a is B = $\dfrac{\mu_0kr^3}4$}
\end{qoptions}
\qn{6} When the number of electrons striking the target material in Coolidge's tube (i.e., X-ray tube) is increased keeping the potential difference same, then :-
\begin{qoptions}
\opt{A}{intensity of X-rays increase}
\opt{B}{wavelength of X-rays photons increase}
\opt{C}{frequency of cut off X-rays photons remains same.}
\opt{D}{energy of most energetic X-rays photons increase.}
\end{qoptions}
\qn{7} A smooth track in the form of a quarter circle of radius 6m lies in the vertical plane. A particle moves from P$_1$ to P$_2$ under the action of forces $\vec F_1,\vec F_2$ and $\vec F_3$. Force $\vec F_1$ is always toward P$_2$ and always 20 N in magnitude. Force $\vec F_2$ is always acts horizontally and is always 30 N in magnitude. Force $\vec F_3$ always acts tangentially to the track and is of magnitude 15 N. Select the correct alternative (s) :\par
\qfig{p1_phy_s2_q07_1}
\begin{qoptions}
\opt{A}{Work done by $\vec F_1$ is 120 J}
\opt{B}{Work done by $\vec F_2$ is 180 J}
\opt{C}{Work done by $\vec F_3$ is 45 $\pi$ J}
\opt{D}{$\vec F_1$ is conservative in nature}
\end{qoptions}
\sect{SECTION-I(iii)}
\qn{8} A solid sphere is rotating about an axis as shown in Figure. An insect follows the dotted path on the circumference of sphere as shown. Match the following :-\par
\qfig{p1_phy_s3_q08_1}
{\footnotesize\setlength{\tabcolsep}{1.5pt}\noindent\begin{tabular}{|c|L{3.2cm}|c|L{4cm}|}
\hline
 & \multicolumn{1}{>{\centering\arraybackslash}p{3.2cm}|}{List$-$I} & & \multicolumn{1}{>{\centering\arraybackslash}p{4cm}|}{List$-$II} \\ \hline
(P) & Moment of inertia & (1) & will remain constant \\ \hline
(Q) & Angular velocity & (2) & will first increase then decrease \\ \hline
(R) & Angular momentum & (3) & will first decrease then increases \\ \hline
(S) & Rotational kinetic energy & (4) & will continuously decrease \\ \hline
 & & (5) & will continuously increase \\ \hline
\end{tabular}}\par
\begin{qoptions}
\opt{A}{P $\to$ 2;Q $\to$ 3;R $\to$ 1;S $\to$ 3}
\opt{B}{P $\to$ 1;Q $\to$ 2;R $\to$ 3;S $\to$ 4}
\opt{C}{P $\to$ 5;Q $\to$ 2;R $\to$ 3;S $\to$ 1}
\opt{D}{P $\to$ 3;Q $\to$ 2;R $\to$ 1;S $\to$ 4}
\end{qoptions}
\qn{9} Match the following\par
{\footnotesize\setlength{\tabcolsep}{1.5pt}\noindent\begin{tabular}{|c|L{1.5cm}|c|L{4.5cm}|}
\hline
 & \multicolumn{1}{>{\centering\arraybackslash}p{1.5cm}|}{List$-$I} & & \multicolumn{1}{>{\centering\arraybackslash}p{4.5cm}|}{List$-$II} \\ \hline
(P) & $\dfrac RL$ & (1) & Time \\ \hline
(Q) & CR & (2) & Frequency \\ \hline
(R) & $\dfrac EB$ & (3) & Speed \\ \hline
(S) & $\dfrac1{\mu_0\varepsilon_0}$ & (4) & Kinetic energy \\ \hline
 & & (5) & Kinetic energy per unit mass \\ \hline
\end{tabular}}\par
\begin{qoptions}
\opt{A}{P $\to$ 3;Q $\to$ 2;R $\to$ 5;S $\to$ 1}
\opt{B}{P $\to$ 2;Q $\to$ 1;R $\to$ 3;S $\to$ 5}
\opt{C}{P $\to$ 1;Q $\to$ 2;R $\to$ 3;S $\to$ 4}
\opt{D}{P $\to$ 5;Q $\to$ 3;R $\to$ 1;S $\to$ 4}
\end{qoptions}
\qn{10} With reference to the following P-V diagram, match the following quantities.\par
[W : work done by gas ; $\Delta$U : change in internal energy of gas]\par
\qfig{p1_phy_s3_q10_1}
{\footnotesize\setlength{\tabcolsep}{1.5pt}\noindent\begin{tabular}{|c|L{2.8cm}|c|L{2cm}|}
\hline
 & \multicolumn{1}{>{\centering\arraybackslash}p{2.8cm}|}{List$-$I} & & \multicolumn{1}{>{\centering\arraybackslash}p{2cm}|}{List$-$II} \\ \hline
(P) & Process MN & (1) & W = 0 \\ \hline
(Q) & Process NO & (2) & W $<$ 0 \\ \hline
(R) & Process OP & (3) & W $>$ 0 \\ \hline
(S) & Process PM & (4) & $\Delta$u $>$ 0 \\ \hline
 & & (5) & $\Delta$u $<$ 0 \\ \hline
\end{tabular}}\par
\begin{qoptions}
\opt{A}{P $\to$ 1,3;Q $\to$ 2,4;R $\to$ 1,5;S $\to$ 1,4}
\opt{B}{P $\to$ 1,4;Q $\to$ 3,4;R $\to$ 2,3;S $\to$ 3,4}
\opt{C}{P $\to$ 1,4;Q $\to$ 3,4;R $\to$ 1,5;S $\to$ 2,5}
\opt{D}{P $\to$ 1,2;Q $\to$ 3,4;R $\to$ 1,5;S $\to$ 2,5}
\end{qoptions}
\qn{11} List$-$I shows different charge distributions and short electric dipole at a distance x from the charge distributions. List$-$II gives the dependence of force acting on the dipole as of function of x.\par
{\footnotesize\setlength{\tabcolsep}{1.5pt}\noindent\begin{tabular}{|c|L{4.6cm}|c|L{1.4cm}|}
\hline
 & \multicolumn{1}{>{\centering\arraybackslash}p{4.6cm}|}{List$-$I} & & \multicolumn{1}{>{\centering\arraybackslash}p{1.4cm}|}{List$-$II} \\ \hline
(P) & \qfigi{p1_phy_s3_q11_P}\par Uniformly charged thin disc. & (1) & x$^{-2}$ \\ \hline
(Q) & \qfigi{p1_phy_s3_q11_Q}\par Uniformly charged thick infinite sheet. Dipole is placed inside the sheet at distance x from axis of sheet. & (2) & x$^{-5}$ \\ \hline
(R) & \qfigi{p1_phy_s3_q11_R}\par Uniform infinite line of charge. & (3) & x$^0$ \\ \hline
(S) & \qfigi{p1_phy_s3_q11_S}\par Uniformly charged sphere. Dipole is placed outside the sphere. & (4) & None of these \\ \hline
 & & (5) & x$^{-3}$ \\ \hline
\end{tabular}}\par
\begin{qoptions}
\opt{A}{P $\to$ 4;Q $\to$ 1;R $\to$ 3;S $\to$ 2}
\opt{B}{P $\to$ 4;Q $\to$ 3;R $\to$ 1;S $\to$ 5}
\opt{C}{P $\to$ 1;Q $\to$ 2;R $\to$ 4;S $\to$ 3}
\opt{D}{P $\to$ 4;Q $\to$ 3;R $\to$ 1;S $\to$ 2}
\end{qoptions}
\sect{SECTION-II}
\qn{1} An ammeter and a voltmeter are connected in series to a battery with an emf E = 6V. When a certain resistance is connected in parallel with the voltmeter, the reading of the latter decreases two times, where as the reading of the ammeter increase the same number of times. Find the voltmeter reading, in volt, after the connection of the resistance.
\qn{2} An ideal gas having initial pressure p, volume V and temperature T is allowed to expand adiabatically until its volume becomes 5.66 V while its temperature falls to $\dfrac T2$. Calculate the degrees of freedom of the gas molecules. Given log$_\mathrm e$(5.66) = 1.73, log$_\mathrm e$ (2) = 0.70
\qn{3} The peak emission from a black body at a certain temperature occurs at a wavelength of 6000 \AA. On increasing its temperature, the total radiation emitted is increased 16 times. These radiations are allowed to fall on a metal surface. Photoelectrons emitted by the peak radiation at higher temperature can be bought to rest by applying a potential equivalent to the excitation potential corresponding to the transition for the level n = 4 to n = 2 in the Bohr's hydrogen atom. The work function of the metal is given by $\dfrac\alpha{100}$eV where $\alpha$ is the numerical constant. The value of $\alpha$ is found 53 $\times$ x, find x. [Take : hc = 12420 eV$-$\AA]
\qn{4} The natural frequency of the system shown in figure is f = $\dfrac1\pi\sqrt{\dfrac{xk}M}$. The pulleys are smooth and massless. Find x.\par
\qfig{p1_phy_s4_q04_1}
\qn{5} A road passes at some distance from a standing man. A truck is coming on the road with some acceleration. The truck driver blows a whistle of frequency 500 Hz when the line joining the truck and the man makes an angle of $\theta$ with the road. The man hears a note having a frequency of 600 Hz when the truck is closest to him. Also, the speed of truck has got doubled during this time. Find the value of $\theta$, in degree.
\qn{6} A uniform disc of radius r and mass m is charged uniformly with the charge q. This disc is placed flat on a rough horizontal surface having coefficient of friction $\mu$. A uniform magnetic field is present in a circular region (a $>$ r) but varying as kt$^3$ as shown in Figure. Find the time, in second after which the disc begins to rotate. (Given r = 1m, m = 18 kg, q = 1 C, $\mu$ = 0.1, K = 4, g = 10 ms$^{-2}$)\par
\qfig{p1_phy_s4_q06_1}
\partsec{2}{CHEMISTRY}
\sect{SECTION-I(i)}
\qn{1} [X] $\mathop{\rightleftharpoons}\limits^{\mathrm{OH}^-}_{\mathrm H^+}$ [Y]\par
Both [X] and [Y] produce ink blue solution with H$_2$O$_2$ in presence of organic solvents and H$_2$SO$_4$, [X] \& [Y] respectively
\begin{qoptions}
\opt{A}{MnO$_4^-$, MnO$_4^{-2}$}
\opt{B}{Cr$_2$O$_7^{-2}$, CrO$_4^{-2}$}
\opt{C}{CrO$_4^{-2}$, Cr$_2$O$_7^{-2}$}
\opt{D}{SO$_3^{-2}$, SO$_4^{-2}$}
\end{qoptions}
\qn{2} Which of the following set contains, all tetrahedral geometries with respect to central metal atom or ion ?
\begin{qoptions}
\opt{A}{FeCl$_4^{2-}$, Cu(CN)$_4^{3-}$, PdCl$_4^{2-}$}
\opt{B}{NiCl$_2$(PPh$_3$)$_2$, Pt(CN)$_4^{2-}$, RhCl(PPh$_3$)$_3$}
\opt{C}{Ni(CO)$_4$, Fe(CO)$_4^{2-}$, Cu(CN)$_4^{3-}$}
\opt{D}{Fe(CO)$_5$ , Co(CO)$_4^-$ , [NiCl$_4$]$^{2-}$}
\end{qoptions}
\qn{3} A small amount of solid lead iodide was added to a beaker of water, which was stirred. Most of the solid settled at the bottom of the beaker, but a little dissolved, establishing the equilibrium\par
PbI$_2$ (s) $\rightleftharpoons$ Pb$^{2+}$ (aq) + 2I$^-$ (aq)\par
The rates of the forward and reverse reactions were monitored over time, producing the graph shown below:\par
\qfig{p1_chem_s1_q03_1}
What happened at time t?
\begin{qoptions}
\opt{A}{The beaker was cooled in an ice-bath.}
\opt{B}{A small amount of solid KI was added to the beaker.}
\opt{C}{A small amount of solid PbI$_2$ was removed from the beaker.}
\opt{D}{A small amount of water was added to the beaker.}
\end{qoptions}
\qfullbegin
\qn{4} Consider the following scheme:\par
\qfig{p1_chem_s1_q04_1}
\qfullsplit
If P is ethyl chloride, R can be\par
\qfig{p1_chem_s1_q04_O}
\begin{qoptions}
\opt{A}{I, II, III, IV}
\opt{B}{II, III, IV only}
\opt{C}{IV only}
\opt{D}{II, IV only}
\end{qoptions}
\qfullend
\sect{SECTION-I(ii)}
\qn{5} One mole of a diatomic ideal gas is subjected to a reversible process that involves two steps AB and AC. The pressure at A and C is same. Consider the graph and choose correct statements.\par
[Given : $\ell$n 2 = 0.70, $\ell$n 3 = 1.1, R = 0.082 L-atm/K-mol]\par
\qfig{p1_chem_s2_q05_1}
\begin{qoptions}
\opt{A}{Entropy change for process AC is 0.2 kJK$^{-1}$}
\opt{B}{In Path AB work done is $-$22.3 J}
\opt{C}{Volume of gas at B = 2 $\times$ Volume of gas at A}
\opt{D}{Volume of gas at C is 32.8 lit.}
\end{qoptions}
\qn{6} The correct statement/s about the pyranose form of a sugar (X) given below is/are:\par
\qfig{p1_chem_s2_q06_1}
\begin{qoptions}
\opt{A}{It exists in two anomeric pyranose forms}
\opt{B}{It reacts with Tollen's reagent to give a silver mirror}
\opt{C}{The penta-O-methyl derivative of (X) is non reducing}
\opt{D}{It resists reduction with aqueous sodium borohydride}
\end{qoptions}
\qn{7} The correct statement(s) related to oxoacids of phosphorous is(are) :
\begin{qoptions}
\opt{A}{Upon heating, H$_3$PO$_3$ undergoes disproportionation reaction to produce H$_3$PO$_4$ and PH$_3$.}
\opt{B}{While H$_3$PO$_3$ can act as reducing agent, H$_3$PO$_4$ cannot.}
\opt{C}{H$_3$PO$_3$ is a monobasic acid.}
\opt{D}{The H atom of P-H bond in H$_3$PO$_3$ is not ionizable in water.}
\end{qoptions}
\sect{SECTION-I(iii)}
\qn{8} Study the sequence of reactions of diborane (B$_2$H$_6$) given below and match the products K, L, M and N with products mention in List-II (shown reagents are taken in excess)\par
\qfig{p1_chem_s3_q08_1}
{\footnotesize\setlength{\tabcolsep}{1.5pt}\noindent\begin{tabular}{|c|L{1cm}|c|L{3.5cm}|}
\hline
 & \multicolumn{1}{>{\centering\arraybackslash}p{1cm}|}{List-I} & & \multicolumn{1}{>{\centering\arraybackslash}p{3.5cm}|}{List-II} \\ \hline
(P) & K & (1) & [BH$_2$(NH$_3$)$_2$][BH$_4$] \\ \hline
(Q) & L & (2) & B$_2$H$_5$Cl \\ \hline
(R) & M & (3) & BCl$_3$ \\ \hline
(S) & N & (4) & B$_3$N$_3$H$_3$Cl$_3$ \\ \hline
 & & (5) & B$_3$N$_3$H$_6$ \\ \hline
\end{tabular}}\par
\begin{qoptions}
\opt{A}{P $\to$ 5;Q $\to$ 1;R $\to$ 3;S $\to$ 2}
\opt{B}{P $\to$ 5;Q $\to$ 1;R $\to$ 2;S $\to$ 4}
\opt{C}{P $\to$ 1;Q $\to$ 5;R $\to$ 2;S $\to$ 4}
\opt{D}{P $\to$ 1;Q $\to$ 5;R $\to$ 3;S $\to$ 4}
\end{qoptions}
\qn{9} Match the half-life of the reaction in List-I with the integrated rate equation in List-II\par
(a = initial concentration of reactant and a $-$ x $-$ concentration of reactant at time `t'.)\par
{\footnotesize\setlength{\tabcolsep}{1.5pt}\noindent\begin{tabular}{|c|L{2.3cm}|c|L{3.6cm}|}
\hline
 & \multicolumn{1}{>{\centering\arraybackslash}p{2.3cm}|}{List $-$ I} & & \multicolumn{1}{>{\centering\arraybackslash}p{3.6cm}|}{List $-$ II} \\ \hline
(P) & $t_{1/2}=\dfrac{0.586}{k}\sqrt a$ & (1) & $k=\dfrac xt$ \\ \hline
(Q) & $t_{1/2}=\dfrac{0.693}{k}$ & (2) & $k=\dfrac1t\left[\dfrac1{a-x}-\dfrac1a\right]$ \\ \hline
(R) & $t_{1/2}=\dfrac a{2k}$ & (3) & $k=\dfrac2t\left[\sqrt a-\sqrt{a-x}\right]$ \\ \hline
(S) & $t_{1/2}=\dfrac1{ka}$ & (4) & $k=\dfrac{2.303}t\log\dfrac a{a-x}$ \\ \hline
 & & (5) & $k=\dfrac2t\left[\dfrac1{\sqrt{a-x}}-\dfrac1{\sqrt a}\right]$ \\ \hline
\end{tabular}}\par
\begin{qoptions}
\opt{A}{P $\to$ 5;Q $\to$ 4;R $\to$ 1;S $\to$ 2}
\opt{B}{P $\to$ 3;Q $\to$ 4;R $\to$ 1;S $\to$ 2}
\opt{C}{P $\to$ 2;Q $\to$ 3;R $\to$ 4;S $\to$ 2}
\opt{D}{P $\to$ 3;Q $\to$ 2;R $\to$ 4;S $\to$ 5}
\end{qoptions}
\qn{10} Match each species in List - I with appropriate statement (s) of List - II and choose correct option\par
{\footnotesize\setlength{\tabcolsep}{1.5pt}\noindent\begin{tabular}{|c|L{3.7cm}|c|L{3cm}|}
\hline
 & \multicolumn{1}{>{\centering\arraybackslash}p{3.7cm}|}{List - I} & & \multicolumn{1}{>{\centering\arraybackslash}p{3cm}|}{List - II} \\ \hline
(P) & \qfigi{p1_chem_s3_q10_P} & (1) & At least one of the product shows optical isomerism. \\ \hline
(Q) & \qfigi{p1_chem_s3_q10_Q} & (2) & Major product contains even number of $\alpha$ - hydrogen's (Hydrogen's which can involve in hyper conjugation). \\ \hline
(R) & \qfigi{p1_chem_s3_q10_R} & (3) & One of the product shows geometrical isomerism. \\ \hline
(S) & \qfigi{p1_chem_s3_q10_S} & (4) & Reactant can turn acidified dichromate to green. \\ \hline
 & & (5) & Will not undergo dehydration. \\ \hline
\end{tabular}}\par
\begin{qoptions}
\opt{A}{P $\to$ 2,3;Q $\to$ 2,4,5;R $\to$ 1,4;S $\to$ 5}
\opt{B}{P $\to$ 2,3,4;Q $\to$ 2,3,4;R $\to$ 1,2;S $\to$ 3,4}
\opt{C}{P $\to$ 1,2,3,4;Q $\to$ 2,3,4;R $\to$ 1,3;S $\to$ 2}
\opt{D}{P $\to$ 2,3,4;Q $\to$ 1,2,3,4;R $\to$ 1,4;S $\to$ 5}
\end{qoptions}
\qn{11} {\footnotesize\setlength{\tabcolsep}{1.5pt}\noindent\begin{tabular}{|c|L{3.6cm}|c|L{2.8cm}|}
\hline
 & \multicolumn{1}{>{\centering\arraybackslash}p{3.6cm}|}{List - I} & & \multicolumn{1}{>{\centering\arraybackslash}p{2.8cm}|}{List - II} \\ \hline
(P) & A compound that liberates ammonia on treatment with hot aqueous KOH & (1) & CH$_3$$-$CH$_2$$-$CN \\ \hline
(Q) & A compound that gives CO$_2$, NH$_4^+$ and acetic acid after boiling in concentrated, aqueous HCl & (2) & CH$_3$$-$CH$_2$$-$CONH$_2$ \\ \hline
(R) & A compound that gives NH$_4^+$ and mono carboxylic acid after treatment with concentrated aqueous HCl. & (3) & NC$-$CH$_2$$-$CN \\ \hline
(S) & A compound that gives product A with MeMgI followed by hydrolysis, which on further treatment with MeMgI excess followed by acidification gives 2-methyl-2-butanol & (4) & CH$_3$COONH$_4$ \\ \hline
 & & (5) & PhCOONH$_4$ \\ \hline
\end{tabular}}\par
The correct option is\_\_\_\_\_\_.
\begin{qoptions}
\opt{A}{P $\to$ 1;Q $\to$ 3;R $\to$ 5;S $\to$ 1}
\opt{B}{P $\to$ 1;Q $\to$ 2;R $\to$ 3;S $\to$ 4}
\opt{C}{P $\to$ 3;Q $\to$ 4;R $\to$ 5;S $\to$ 4}
\opt{D}{P $\to$ 2;Q $\to$ 4;R $\to$ 2;S $\to$ 1}
\end{qoptions}
\sect{SECTION-II}
\qn{1} Find the number of reactions, which give N$_2$.\par
(I) (NH$_4$)$_2$Cr$_2$O$_7\xrightarrow{\Delta}$ \quad (II) Ba(N$_3$)$_2\xrightarrow{\Delta}$\par
(III) NH$_4$Cl $\xrightarrow{\Delta}$ \quad (IV) NH$_3$ (excess) + Cl$_2$ $\rightarrow$\par
(V) NH$_4$NO$_3$ $\xrightarrow[250^\circ\mathrm C]{\Delta}$ \quad (VI) (NH$_4$)$_2$SO$_4$ $\xrightarrow[\mathrm{aq.}]{\mathrm{KOH}}$\par
(VII) C$_2$H$_5$NH$_2$ + HNO$_2$ $\xrightarrow{\mathrm{water}}$
\qn{2} If `x' is the total number of p$\pi$ $-$d$\pi$ bonds and `y' is the total number of lone pair(s) on non terminal atoms per molecule in trimer of SO$_3$, the value of (x + y) is :-
\qn{3} One gram of commercial AgNO$_3$ is dissolved in 50 mL of water. It is treated with 50 mL of a KI solution. The silver iodide thus precipitated is filtered off. Excess of KI in the filtrate is titrated with M/10 KIO$_3$ solution in presence of 6 M HCl till all I$^-$ ions are converted into ICl. It requires 50 mL of M/10 KIO$_3$ solution. Twenty milliliters of the same stock solution of KI requires 30 mL of M/10 KIO$_3$ under similar conditions. The percentage of AgNO$_3$ in the sample is (Atomic weight of silver = 108)
\qn{4} Both gold and silver are extracted by cyanide process. Gold in nature is frequently alloyed with silver, which is also oxidized by aerated NaCN solution, similar to gold. 1000 L of a solution of 0.002 M [Au(CN)$_2$]$^-$ and 0.003 M [Ag(CN)$_2$]$^-$ was evaporated to half of the original volume and treated with 78 g zinc (atomic mass = 65).\par
Assume that elevation from standard conditions is unimportant and that the redox reactions go to completion. Given:\par
[Zn(CN)$_4$]$^{2-}$ + 2e$^-$ $\to$ Zn (s) + 4CN$^-$; E$^\circ$ = $-$1.26V\par
[Au(CN)$_2$]$^-$ + e$^-$ $\to$ Au (s) + 2CN$^-$; E$^\circ$ = $-$0.60V\par
[Ag(CN)$_2$]$^-$ + e$^-$ $\to$ Ag (s) + 2CN$^-$; E$^\circ$ = $-$0.31V\par
If the concentration of [Au(CN)$_2$]$^-$in the solution after addition of zinc in it, be x mol/lit. The value of (x $\times$ 10$^5$) is.
\qn{5} Consider the following reactions\par
\qfig{p1_chem_s4_q05_1}
Amount (in mg) of B reacts completely with 100 ml of 0.1 M NaHCO$_3$ solution.
\qn{6} The correct order of the given reagents to convert benzene to m-Chlorobromobenzene is\par
(1) Sulfuric acid (conc.) and heat\par
(2) Cl$_2$ + FeCl$_3$ and heat\par
(3) NaNO$_3$ + H$_3$O$^+$, 0$^\circ$C\par
(4) H$_2$/Pt catalyst\par
(5) Mg in ether\par
(6) PBr$_3$\par
(7) H$_3$PO$_2$ (aqueous)\par
(8) HNO$_3$ (conc.) + H$_2$SO$_4$ (conc.) and heat\par
(9) Cu$_2$Br$_2$ +HBr\par
(10) (CH$_3$CO)$_2$ O + pyridine\par
Write the correct order in form of their number and express last two digits as answer.\par
[ For, e.g. if order is 1357, then your answer will be 57]
\partsec{3}{MATHEMATICS}
\sect{SECTION-I(i)}
\qn{1} A ray of light travelling along the line OP (O being origin) is reflected by the line mirror 2x $-$ 3y + 1 = 0, the point of incidence being P(1,1). The reflected ray, travelling along PQ is again reflected by the line mirror 2x $-$ 3y $-$ 1 = 0, the point of incidence being Q. From Q, ray move along QR, where R lies on the line 2x $-$ 3y + 1 = 0. Then sum of ordinate of point R and abscissa of point Q is :
\begin{qoptions}
\opt{A}{$\dfrac{100}{13}$}
\opt{B}{$\dfrac{73}{13}$}
\opt{C}{$\dfrac{53}{13}$}
\opt{D}{$\dfrac{23}{13}$}
\end{qoptions}
\qn{2} Number of integers in the range of `m' for which point (m,1) lies in the smaller region bounded by x$^2$ + y$^2$ $-$ 3x + 1 = 0 and 2x $-$ y = 2 is-
\begin{qoptions}
\opt{A}{0}
\opt{B}{1}
\opt{C}{2}
\opt{D}{3}
\end{qoptions}
\qn{3} Let $g(x)=\displaystyle\int_{x^2}^{x^3}\dfrac{dt}{\ln t}$ for x $>$ 1 and $h(x)=\displaystyle\int_1^x\left(2t^2-\ln t\right)g(t)dt$ (x $>$ 1), then :
\begin{qoptions}
\opt{A}{h is increasing on (1, $\infty$)}
\opt{B}{h is decreasing on (1, 2) and increasing on (2, $\infty$)}
\opt{C}{h is decreasing on (1, $\infty$)}
\opt{D}{h is increasing on (1, 2) and decreasing on (2, $\infty$)}
\end{qoptions}
\qn{4} The tangent at (x$'$, y$'$) to the curve x$^3$ + y$^3$ = a$^3$ meets the curve again at (x$_1$, y$_1$), then $\left|\dfrac{x_1}{x'}+\dfrac{y_1}{y'}\right|^2=$
\begin{qoptions}
\opt{A}{4}
\opt{B}{1}
\opt{C}{a}
\opt{D}{a$^2$}
\end{qoptions}
\sect{SECTION-I(ii)}
\qn{5} Let f(x) = sin x + px + q. Then f(x) = 0 has :
\begin{qoptions}
\opt{A}{Only one real root which is positive if p $>$ 1, q $<$ 0}
\opt{B}{Only one real root which is negative if p $<$ $-$1, q $<$ 0}
\opt{C}{Only one real root which is negative if p $>$ 1, q $>$ 0}
\opt{D}{Only one real root which is positive if p $<$ $-$ 1, q $<$ 0}
\end{qoptions}
\qn{6} Let $\vec u=\cos\alpha\hat i+\sin\alpha\hat j$, $\vec v=\cos\beta\hat i+\sin\beta\hat j$, $\vec w=\cos\gamma\hat i+\sin\gamma\hat j$, $0\leq\alpha,\beta,\gamma<2\pi$. If $\vec u.\vec w=\vec v.\vec w=-\sqrt3\vec u.\vec v$ then $\left(2\vec u.\vec v\right)^3$ is equal to :
\begin{qoptions}
\opt{A}{$\dfrac{-8}{27}$}
\opt{B}{$\dfrac{-1}8$}
\opt{C}{1}
\opt{D}{$\dfrac19$}
\end{qoptions}
\qn{7} P is a non$-$singular matrix of order n $\times$ n such that 3PQP$^{-1}$ + P = 2P$^{-1}$ QP then
\begin{qoptions}
\opt{A}{P \& Q both are identity matrices}
\opt{B}{det.(P + Q) = 0}
\opt{C}{PQP$^{-1}$ $-$ P$^{-1}$QP is a singular matrix}
\opt{D}{P + Q is non - invertible matrix}
\end{qoptions}
\sect{SECTION-I(iii)}
\qn{8} Matching the following statement\par
{\footnotesize\setlength{\tabcolsep}{1.5pt}\noindent\begin{tabular}{|c|L{5.2cm}|c|L{0.9cm}|}
\hline
 & \multicolumn{1}{>{\centering\arraybackslash}p{5.2cm}|}{List-I} & & \multicolumn{1}{>{\centering\arraybackslash}p{0.9cm}|}{List-II} \\ \hline
(P) & If $\dfrac{dy}{dx}\left(\dfrac{1+\cos x}y\right)=-\sin x$ and f$\left(\dfrac\pi2\right)=-1$ then f(0) is & (1) & 5 \\ \hline
(Q) & The number of points of non-derivability of the function f(x) = $\left[\dfrac{2x}\pi\right]$ sgn$\left(\dfrac1{\{x\}}\right)$ in ($-$2, 2) is (where [P], \{P\} \& sgn(P) denotes greater integer $\leq$ P, fractional part of P \& signum function of P respectively) & (2) & $-$2 \\ \hline
(R) & The area bounded by $e^{\ln|x+1|}\geqslant|y|,|x|\leqslant1$ is & (3) & 4 \\ \hline
(S) & Number of real values of x which satisfy the equation $\dfrac{\left(\sec^{-1}x\right)^3+\left(\mathrm{cosec}^{-1}x\right)^3}{\left(\tan^{-1}x+\cot^{-1}x\right)^3}=7$ is: & (4) & 1 \\ \hline
 & & (5) & 7 \\ \hline
\end{tabular}}\par
\begin{qoptions}
\opt{A}{P $\to$ 1;Q $\to$ 2;R $\to$ 3;S $\to$ 4}
\opt{B}{P $\to$ 2;Q $\to$ 1;R $\to$ 3;S $\to$ 4}
\opt{C}{P $\to$ 2;Q $\to$ 3;R $\to$ 1;S $\to$ 4}
\opt{D}{P $\to$ 2;Q $\to$ 1;R $\to$ 4;S $\to$ 3}
\end{qoptions}
\qn{9} Matching the following statement\par
{\footnotesize\setlength{\tabcolsep}{1.5pt}\noindent\begin{tabular}{|c|L{5.2cm}|c|L{0.9cm}|}
\hline
 & \multicolumn{1}{>{\centering\arraybackslash}p{5.2cm}|}{List-I} & & \multicolumn{1}{>{\centering\arraybackslash}p{0.9cm}|}{List-II} \\ \hline
(P) & If sin$^2$ A = x and sinA $\cdot$ sin2A sin3A $\cdot$ sin4A = ax$^2$ + bx$^3$ + cx$^4$ + dx$^5$, then the value of a + b + c + d + 5 is & (1) & 1 \\ \hline
(Q) & The values of x between 0 and 2$\pi$ which satisfy the equation $\sin x\sqrt{8\cos^2x}=1$ are in A.P. with common difference d, then $\dfrac{4d}\pi$ is & (2) & 2 \\ \hline
(R) & The least difference between the roots, in the first quadrant $\left(0\leqslant x\leqslant\dfrac\pi2\right)$, of the equation $4\cos x\left(2-3\sin^2x\right)+(\cos2x+1)=0$ is $\dfrac\pi k$, then k = & (3) & 6 \\ \hline
(S) & Let $f(x)=\sqrt{\sin^4x+4\cos^2x}-\sqrt{\cos^4x+4\sin^2x}$ and g(sin 2t) = sint + cos t $\forall$ t $\in\left[\dfrac{-\pi}4,\dfrac\pi4\right]$. If the range of g(f(x)) is [a, b], then the value of a$^2$ + b$^2$ is\_\_\_\_\_\_. & (4) & 5 \\ \hline
 & & (5) & 4 \\ \hline
\end{tabular}}\par
\begin{qoptions}
\opt{A}{P $\to$ 2;Q $\to$ 4;R $\to$ 1;S $\to$ 3}
\opt{B}{P $\to$ 4;Q $\to$ 1;R $\to$ 2;S $\to$ 3}
\opt{C}{P $\to$ 2;Q $\to$ 1;R $\to$ 4;S $\to$ 3}
\opt{D}{P $\to$ 4;Q $\to$ 1;R $\to$ 3;S $\to$ 2}
\end{qoptions}
\qn{10} Matching the following statement\par
{\footnotesize\setlength{\tabcolsep}{1.5pt}\noindent\begin{tabular}{|c|L{5.2cm}|c|L{0.9cm}|}
\hline
 & \multicolumn{1}{>{\centering\arraybackslash}p{5.2cm}|}{List- I} & & \multicolumn{1}{>{\centering\arraybackslash}p{0.9cm}|}{List-II} \\ \hline
(P) & If $\displaystyle\lim_{h\to0^+}\int_{-1}^1\dfrac{h\,dx}{h^2+x^2}=k\pi$ Then k equals & (1) & 5 \\ \hline
(Q) & The area of the region defined by 1 $\leqslant$ $|x-2|+|y+1|$ $\leqslant$ 2 is & (2) & 3 \\ \hline
(R) & $\displaystyle\int_0^{110\pi}[\sin x+\cos x]\,dx=-11k\pi$ then k equal to (where [.] denotes greatest integer function) & (3) & 1 \\ \hline
(S) & If distance between two focii of ellipse $(x-2y+3)^2+(8x+4y+4)^2=20$ is k then $\dfrac{k^2}5$ equals to & (4) & 6 \\ \hline
 & & (5) & 9 \\ \hline
\end{tabular}}\par
\begin{qoptions}
\opt{A}{P $\to$ 3;Q $\to$ 4;R $\to$ 1;S $\to$ 2}
\opt{B}{P $\to$ 2;Q $\to$ 3;R $\to$ 1;S $\to$ 4}
\opt{C}{P $\to$ 3;Q $\to$ 4;R $\to$ 2;S $\to$ 1}
\opt{D}{P $\to$ 2;Q $\to$ 4;R $\to$ 1;S $\to$ 3}
\end{qoptions}
\qn{11} {\footnotesize\setlength{\tabcolsep}{1.5pt}\noindent\begin{tabular}{|c|L{5.2cm}|c|L{0.9cm}|}
\hline
 & \multicolumn{1}{>{\centering\arraybackslash}p{5.2cm}|}{List- I} & & \multicolumn{1}{>{\centering\arraybackslash}p{0.9cm}|}{List-II} \\ \hline
(P) & If $\sqrt{x^2+y^2}+\sqrt{(x+3)^2+(y+4)^2}=9-a^2$ represent an ellipse then number of integral value of a is & (1) & 7 \\ \hline
(Q) & If a, b, c $>$ 0 \& a$^2$ = bc \& a + b + c = abc then the least possible value of a$^2$ is & (2) & 5 \\ \hline
(R) & The minimum value of n $\in$ N so that the sum of coefficient of terms in the expansion of (a +b + c $-$ d)$^n$ is more than the number of terms is:- & (3) & 3 \\ \hline
(S) & If $\displaystyle\lim_{x\to\infty}\left(8x^3+ax^2\right)^{\frac13}-bx$ exists \& is equal to 1, then value of a $-$ 3b is equals to & (4) & 6 \\ \hline
 & & (5) & 2 \\ \hline
\end{tabular}}\par
\begin{qoptions}
\opt{A}{P $\to$ 3;Q $\to$ 4;R $\to$ 1;S $\to$ 1}
\opt{B}{P $\to$ 2;Q $\to$ 1;R $\to$ 4;S $\to$ 1}
\opt{C}{P $\to$ 2;Q $\to$ 1;R $\to$ 4;S $\to$ 4}
\opt{D}{P $\to$ 3;Q $\to$ 3;R $\to$ 1;S $\to$ 4}
\end{qoptions}
\sect{SECTION-II}
\qn{1} The solution of the equation $\left(2x^2+3y^2-7\right)x\,dx-\left(3x^2+2y^2-8\right)y\,dy=0$ is given by $\left(x^2+y^2+\gamma\right)=C\left(x^2-y^2+\delta\right)^\beta$. Where C is arbitrary constant, then ($\beta$ + $\gamma$ + $\delta$) equals
\qn{2} If z is a complex number satisfying the condition $\left|z-\dfrac5z\right|=4$, then sum of greatest and least value of $|z|$ is
\qn{3} If T$_\mathrm n$ is term which is independent of `x' in the expansion of $\left(\dfrac{x^4+x^2+1}{x^2-x+1}-\dfrac{x^2+x+1}x\right)^{12}$, then n is
\qn{4} In a drawer Mr. X has 5 pairs of socks, each pair a different color. On Monday, he selects two individual socks at random from the 10 socks in the drawer. On Tuesday, he selects 2 of the remaining 8 socks at random and on Wednesday, two of the remaining 6 socks at random. If the probability that Wednesday is the first day he selects matching socks is $\dfrac mn$, where m and n are relatively prime positive integers, then the value of $\dfrac{m+n}{11}$ is equal to
\qn{5} A 3 $\times$ 3 determinant has entries either 1 or $-$1. Let S be the set of all determinants such that the product of elements of any row or column is $-$1. For example $\begin{vmatrix}1&-1&1\\1&1&-1\\-1&1&1\end{vmatrix}$ is an element of S and number of elements of S is m. Let P = $\begin{vmatrix}3&-2&3\\2&-2&3\\0&-1&1\end{vmatrix}$ and trace of the matrix adj (adj P) is n then the value of $\dfrac mn$ is
\qn{6} Let the area of given triangle ABC be $\Delta$. Points A$_1$, B$_1$ and C$_1$ are the mid points of the BC, CA and AB respectively. Point A$_2$ is the mid point of CA$_1$. Lines C$_1$A$_1$ and AA$_2$ meet the median BB$_1$ at points E and D respectively. If $\Delta_1$ be the area of the quadrilateral A$_1$A$_2$DE. If the value of $\dfrac{\Delta_1}\Delta=\dfrac mn$ (where m and n are coprime), then n $-$ 3m = \_\_\_\_\_.
\raggedcolumns
\papersec{2}
\partsec{1}{PHYSICS}
\sect{SECTION-I(i)}
\qn{1} A current carrying wire is placed along ACB in the grooves of an insulating semicircular disc of radius `R' as shown. The current enters at point A and leaves at point B. Determine the magnetic field at point D.\par
\qfig{p2_phy_s1_q01_1}
\begin{qoptions}
\opt{A}{$\dfrac{\mu_0i}{8\pi R\sqrt3}$}
\opt{B}{$\dfrac{\mu_0i}{4\pi R\sqrt3}$}
\opt{C}{$\dfrac{\sqrt3\mu_0i}{4\pi R}$}
\opt{D}{None}
\end{qoptions}
\qn{2} In the following radioactive decay C is stable nucleus. Then\par
$\mathrm A\xrightarrow{\lambda_1}\mathrm B\xrightarrow{\lambda_2}\mathrm C$\par
time t = 0 \quad N$_0$ \quad 0 \quad 0\par
time t \quad N$_1$ \quad N$_2$ \quad N$_3$
\begin{qoptions}
\opt{A}{Number of nuclei of B will first increase and then decrease.}
\opt{B}{Rate of decay of A will first increase and then decrease.}
\opt{C}{If $\lambda_2>\lambda_1$ then activity of B will always be higher than activity of A.}
\opt{D}{If $\lambda_1>\lambda_2$ then number of nuclei of C will always be less than number of nuclei of B.}
\end{qoptions}
\qn{3} In the given figure mass of block A is m and that of B is 2m. They are attached at the ends of a spring and kept on a smooth horizontal surface. Now, the spring is compressed by x$_0$ and released. Find displacement of block B by the time compression of the spring is reduced to $\dfrac{x_0}2$ :-\par
\qfig{p2_phy_s1_q03_1}
\begin{qoptions}
\opt{A}{$\left(\dfrac{x_0}3\right)$}
\opt{B}{$\left(\dfrac{x_0}2\right)$}
\opt{C}{$\left(\dfrac{x_0}6\right)$}
\opt{D}{$\left(\dfrac{x_0}4\right)$}
\end{qoptions}
\qn{4} A rectangular tank is filled completely with water. A hole at its bottom is unplugged. The graph between the velocity of efflux (through a small hole) vs depth of water h from the top of tank is :-\par
\qfig{p2_phy_s1_q04_1}
\begin{qoptions}
\opt{A}{\qfigi{p2_phy_s1_q04_A}}
\opt{B}{\qfigi{p2_phy_s1_q04_B}}
\opt{C}{\qfigi{p2_phy_s1_q04_C}}
\opt{D}{\qfigi{p2_phy_s1_q04_D}}
\end{qoptions}
\sect{SECTION-I(ii)}
\qn{5} The main scale of a vernier calipers reads in mm and its vernier is divided into 10 divisions which coincides with 9 divisions of the main scale. It was also observed when two jaws are brought in contact the zero of the vernier scale coincided with the zero of the main scale. The length of a cube is measured with this vernier calipers. The main scale reads 10 mm and the first division of the vernier scale coincides with the main scale. If the mass of the cube is 2.0606 g then mark the correct alternatives.
\begin{qoptions}
\opt{A}{The density of the cube in appropriate significant figures is 2.0000 gm/cm$^3$.}
\opt{B}{The density of the cube in appropriate significant figures is 2.00 g/cm$^3$.}
\opt{C}{If the zero mark of the vernier scale lies towards the left side of the zero of the main scale when the two jaws are brought in contact, the error will be positive and the correction will be negative.}
\opt{D}{If the zero mark of the vernier scale lies towards the right side of the zero of the main scale when two jaws are brought in contact the error will be positive and the correction will be negative.}
\end{qoptions}
\qn{6} For the given potential versus distance graph for one dimensional electric field and a particle of charge 20 $\mu$C and mass 0.2 kg, Choose the correct option(s).\par
\qfig{p2_phy_s2_q06_1}
\begin{qoptions}
\opt{A}{Force acting on the particle at x = 25 cm is zero.}
\opt{B}{Force acting on the particle at x = 25 cm is 4$\hat{\mathrm i}$ Newton.}
\opt{C}{Minimum speed at origin required by the particle to reach x = 25 cm is 2 m/s.}
\opt{D}{Minimum speed at origin required by the particle to reach x = 25 cm is $\sqrt2$ m/s.}
\end{qoptions}
\qn{7} A sky lab of mass m is first launched from the surface of earth (radius R) in a circular orbit of radius 2R (from centre of earth) and then it is shifted from this circular orbit to another circular orbit of radius 3R.
\begin{qoptions}
\opt{A}{Minimum energy required to place the lab in first orbit is equal to $\dfrac34$ mgR.}
\opt{B}{Minimum energy required to shift the lab from first orbit to second orbit is $\dfrac{mgR}{12}$.}
\opt{C}{Energy of sky lab in second orbit is greater than its energy in first orbit.}
\opt{D}{Energy of sky lab in second orbit is less than its energy in first orbit.}
\end{qoptions}
\qn{8} P-V diagram of cyclic process ABCA is as shown in figure. Choose the correct statement(s)\par
\qfig{p2_phy_s2_q08_1}
\begin{qoptions}
\opt{A}{$\Delta\mathrm Q_{\mathrm A\to\mathrm B}$ = +ve}
\opt{B}{$\Delta\mathrm U_{\mathrm B\to\mathrm C}$ = +ve}
\opt{C}{$\Delta\mathrm U_{\mathrm{CA}}$ = $-$ve}
\opt{D}{$\Delta\mathrm W_{\mathrm{CAB}}$ = $-$ve}
\end{qoptions}
\qn{9} A ball of mass m moving horizontally at a speed v collides with the bob of a simple pendulum at rest. The mass of the bob is also m. Here v $\leqslant$ $\sqrt{2g\ell}$. (Where $\ell$ is length of the pendulum)
\begin{qoptions}
\opt{A}{If the collision is perfectly inelastic, the height to which the two balls rise after the collision is $\dfrac{v^2}{8g}$.}
\opt{B}{If the collision is perfectly inelastic, the kinetic energy of the system immediately after the collision becomes half of that before collision.}
\opt{C}{If the collision is perfectly elastic, the bob of the pendulum will rise to a height of $\dfrac{v^2}{2g}$.}
\opt{D}{If the collision is perfectly elastic, the kinetic energy of the system immediately after the collision is equal to that before collision.}
\end{qoptions}
\qn{10} A pulley of radius R = 20 cm is rotated by a weight suspended from a thread which is gradually being wound off the pulley. Initially the weight is at rest and then it begins to go down with an acceleration a = 2 cm/s$^2$ :-\par
\qfig{p2_phy_s2_q10_1}
\begin{qoptions}
\opt{A}{Angular velocity of the pulley when the weight has travelled a distance S$_1$ = 100 cm is 1rad/s.}
\opt{B}{Angular velocity of the pulley when the weight has travelled a distance S$_1$ = 100 cm is 2rad/s.}
\opt{C}{The magnitude of the acceleration of point A when the weight has travelled a distance S$_1$ = 100 cm is $20\sqrt2$ cm/s$^2$.}
\opt{D}{The magnitude of the acceleration of point A when the weight has travelled a distance S$_1$ = 100 cm is 20.1 cm/s$^2$.}
\end{qoptions}
\sect{SECTION-III}
\qn{1} Two thin lenses with lens powers D$_1$ and D$_2$ diopter are located at distance L = 25 cm from each other and their optical axes coincide. This system creates an erect real image of the object, located at the optical axis closer to lens D$_1$ with the net magnification M$_1$ = 1. If the position of the two lenses are exchanged, the system again produces an erect real image with the net magnification M$_2$ = 4. What is the difference between lens powers $\Delta$D = D$_1$ $-$ D$_2$ (in diopter)?
\qn{2} A storage battery (E = 24 V) was connected to one end of a two conductor telephone line of length L=5.6 km. There is a leakage somewhere in the line. If the free ends are open, then current through the battery is I$_1$ = 1.6 A and if the free ends of the line are short-circuited then current through the battery is I$_2$ = 2 A. The short-circuit current of the battery is I$_3$ = 24 A and resistance of each conductor is r = 7$\Omega$. Determine the resistance (in $\Omega$) of the insulation R at the point of the breakdown.
\qn{3} A string will break apart if it is placed under too much tensile stress. One type of steel has density $\rho_{\mathrm{steel}}$ = 10$^4$ kg/m$^3$ and breaking stress $\sigma$ = 9 $\times$ 10$^8$ N/m$^2$. We make a guitar string from (4$\pi$) gram of this type of steel. It should be able to withstand (900$\pi$) N without breaking. The highest possible fundamental frequency (in Hz) of standing waves on the string is 75n. Find n.
\qn{4} A car begins to move at time t = 0 and then accelerates along a straight track with a velocity given by V(t) = 2t$^2$ ms$^{-1}$ for 0 $\leq$ t $\leq$ 2, where t is time in second. After the end of acceleration, the car continues to move at a constant speed. A small block initially at rest on the floor of the car begins to slip at t = 1 s and stops slipping at t = 3 s. The coefficient of static and kinetic friction between the block and the floor are $\mu_\mathrm s$ and $\mu_\mathrm k$ respectively. Find the value of $\dfrac{3\mu_\mathrm s}{\mu_\mathrm k}$.
\qn{5} A plane spiral with a great number N of turns wound tightly to one another is located in a uniform magnetic field perpendicular to the spiral's plane. The outside radius of the spiral's turns is equal to a and inner radius is zero. The magnetic induction varies with time as B = B$_0$ sin$\omega$t, where B$_0$ and $\omega$ are constants. Find the amplitude of emf induced in the spiral. (Given : $\omega=\dfrac1\pi$ rad/s, N = 1000, a = 3cm, B$_0$ = 10 Tesla)
\qn{6} A steel rod is 4.000 cm in diameter at 30$^\circ$C. A brass ring has an interior diameter of 3.992 cm at 30$^\circ$C. In order that the ring just slides onto the steel rod, the common temperature of the two should be $\theta$ (in $^\circ$C). Find $\theta$/100 rounded off to nearest integer ($\alpha_{\mathrm{steel}}$ = 11 $\times$ 10$^{-6}$/$^\circ$C and $\alpha_{\mathrm{brass}}$ = 19 $\times$ 10$^{-6}$/$^\circ$C)
\qn{7} A particle of mass 1 mg and charge 1$\mu$C is projected with speed 50 m/s towards a region containing a uniform magnetic field 10 T. What should be the minimum thickness (d) of this region (in m) so that the particle returns back to left side?\par
\qfig{p2_phy_s3_q06_1}
\qn{8} A point charge is moving along a straight line with constant velocity v. Consider a small area A perpendicular to the direction of motion of charge as shown in figure.\par
\qfig{p2_phy_s3_q08_1}
The displacement current through the area when its distance from the charge is x is $\dfrac{qAv}{2\pi x^n}$. The value of n is :-
\partsec{2}{CHEMISTRY}
\sect{SECTION-I(i)}
\qn{1}\par
\qfig{p2_chem_s1_q01_1}
The correct statement regarding product P \& Q is :
\begin{qoptions}
\opt{A}{Both are optically active esters}
\opt{B}{P \& Q are diastereomers}
\opt{C}{P \& Q are meso compounds}
\opt{D}{Out of P \& Q, one is optically active and other is optically inactive.}
\end{qoptions}
\qn{2} Inversion of sucrose (C$_{12}$H$_{22}$O$_{11}$) is first-order reaction and is studied by measuring angle of rotation of the reaction mixture at different instants of time\par
C$_{12}$H$_{22}$O$_{11}$ + H$_2$O $\xrightarrow{\mathrm H^+}$ C$_6$H$_{12}$O$_6$ + C$_6$H$_{12}$O$_6$\par
If r$_0$, r$_\mathrm t$ and r$_\infty$ are the angles of rotation of the reaction mixture at the start, at the time t and at the end of the reaction respectively, then 50\% inversion of sucrose occurs when :
\begin{qoptions}
\opt{A}{r$_0$ = 2r$_\mathrm t$ $-$ r$_\infty$}
\opt{B}{r$_0$ = r$_\mathrm t$ $-$ r$_\infty$}
\opt{C}{r$_0$ = r$_\mathrm t$ $-$ 2r$_\infty$}
\opt{D}{r$_0$ = r$_\mathrm t$ + r$_\infty$}
\end{qoptions}
\qn{3} CoS (Black) $\xrightarrow{\text{aque--regia}}$ solution $\xrightarrow[\mathrm{Br_2\ water}]{\mathrm{exc.NaHCO_3}}$ coloured complex (P)\par
What is co-ordination no. of central metal in complex (P)
\begin{qoptions}
\opt{A}{5}
\opt{B}{6}
\opt{C}{3}
\opt{D}{4}
\end{qoptions}
\qn{4} What is general formual of below mentioned silicate ? [o : oxygen; $\bullet$ : silicon]\par
\qfig{p2_chem_s1_q04_1}
\begin{qoptions}
\opt{A}{$\left(\mathrm{Si_4O_{11}^{6-}}\right)_n$}
\opt{B}{$\left(\mathrm{Si_2O_5^{2-}}\right)_n$}
\opt{C}{$\left(\mathrm{Si_6O_{17}^{10-}}\right)_n$}
\opt{D}{$\left(\mathrm{Si_3O_8^{4-}}\right)_n$}
\end{qoptions}
\sect{SECTION-I(ii)}
\qn{5}\par
\qfig{p2_chem_s2_q05_S}
Product Y and Z are separated and then heated with NaOH. Y gave P where as Z gave Q and CH$_3$OH as product then which of the following is/are correct :
\begin{qoptions}
\opt{A}{Product P is \qfigi{p2_chem_s2_q05_P}}
\opt{B}{Product Z is $\mathrm{H-\overset{O}{\overset{\|}{C}}-H}$}
\opt{C}{Product Y is \qfigi{p2_chem_s2_q05_Y}}
\opt{D}{Product X is \qfigi{p2_chem_s2_q05_X}}
\end{qoptions}
\qn{6} Which of the following is/are resolvable ?
\begin{qoptions}
\opt{A}{\qfigi{p2_chem_s2_q06_a}}
\opt{B}{\qfigi{p2_chem_s2_q06_b}}
\opt{C}{\qfigi{p2_chem_s2_q06_c}}
\opt{D}{\qfigi{p2_chem_s2_q06_d}}
\end{qoptions}
\qn{7} Select the correct statement(s) from the following.
\begin{qoptions}
\opt{A}{If the critical micelle concentration (CMC) of a soap is 10$^{-3}$ M, then 10$^{-4}$ M solution of this soap will have colloidal nature}
\opt{B}{Charge on the colloidal particles may be determined with the help of dialysis.}
\opt{C}{Peptization is the process of conversion of a colloidal sol into precipitate by adding some electrolyte.}
\opt{D}{When some oil-soluble dye is added in the water in oil emulsion, the background becomes coloured.}
\end{qoptions}
\qn{8} EMF of the cell: Cd(s) $|$ CdCl$_2$.5H$_2$O (sat.) $|$ AgCl(s) $|$ Ag(s) is +0.70 V at 0$^\circ$C and +0.60 V at 50$^\circ$C. If $\Delta$H$^\circ$ and $\Delta$S$^\circ$ are temperature independent, then the correct information(s) regarding the cell reaction is/are
\begin{qoptions}
\opt{A}{$\Delta$G$^\circ$ = $-$115.8 kJ at 50$^\circ$C}
\opt{B}{$\Delta$G$^\circ$ = 135.1 kJ at 0$^\circ$C}
\opt{C}{$\Delta$S$^\circ$ = $-$386 J/K}
\opt{D}{$\Delta$H$^\circ$ = $-$221.178 kJ}
\end{qoptions}
\qn{9} Select the correct comparison(s) of properties indicated in options :
\begin{qoptions}
\opt{A}{MnO$_4^{2-}$ $>$ MnO$_4^-$ : (paramagnetic character)}
\opt{B}{Na$_2$Cr$_2$O$_7$ $>$ K$_2$Cr$_2$O$_7$ (Solubility in water)}
\opt{C}{Lu$^{+3}$ $<$ Ce$^{+3}$ $<$ La$^{+3}$ : (Ionic radius)}
\opt{D}{Mn$^{+3}$ $<$ Fe$^{+3}$ $<$ Ti$^{+3}$ : (oxidising power)}
\end{qoptions}
\qn{10} The extraction of pure metal from bauxite ore (containing iron oxide as chief impurity) involves :
\begin{qoptions}
\opt{A}{Electrolytic reduction}
\opt{B}{Poling}
\opt{C}{Leaching by exc. NaOH(aq).}
\opt{D}{Precipitation of amphoteric hydroxide}
\end{qoptions}
\sect{SECTION-III}
\qn{1} A container of capacity 20 L contains N$_2$(g), water vapour and negligible but sufficient amount of liquid water to maintain equilibrium, at 300K and 1 atm. The liquid water is now decomposed completely into H$_2$(g) and O$_2$(g), by any means, without changing the temperature. If the final pressure becomes 1.9 atm, the mass of liquid water present initially (in gm) was (Nearest integer)\par
(Aqueous tension of H$_2$O at 300K = 0.04 atm, R = 0.08 L-atm/K-mole)
\qn{2} Cd$^{2+}$ ion is precipitate after passing H$_2$S in acidified solution. In the final solution 0.1 milli-mole of CdSO$_4$ are present in 10 ml acidic saturated solution of 0.10M H$_2$S. What is the pH of the solution after filtering off precipitate, boiling off H$_2$S and making the solution 100 ml by adding H$_2$O :\par
[K$_\mathrm{sp}$ of CdS = 10$^{-28}$, For H$_2$S : K$_{\mathrm a_1}$ = 10$^{-9}$, K$_{\mathrm a_2}$ = 10$^{-20}$]
\qn{3} Let us assume that de$-$Broglie equation is valid in Bohr model. In H$-$like sample, an electron is in an orbit which has path length of the orbit is equal to 3 times of wavelength of electrons in that orbit. This orbit of H$-$like sample has total energy equal to the ground state energy of H$-$atom. Now we have same sample having all electrons in the ground state. If photon of energy 126.566 eV is given to this ground state sample then approximate value of wavelength ($\lambda$)(in \AA, closest integer) of ejected e$^-$ is :
\qn{4} Iodine reacts with chlorine in presence of water to give an iodine containing compound `x'. How many moles of `x' is formed on reaction of 2 moles I$_2$ with 10 moles of Cl$_2$
\qn{5} A solution contains equimolar amount of following ions : Zn$^{+2}$, Fe$^{+3}$, Ag$^+$, Cu$^{+2}$, Ni$^{+2}$, Pb$^{+2}$, Cd$^{+2}$, Hg$^{+2}$, Cr$^{+3}$. How many of cations form soluble complex with excess of aqueous NH$_3$?
\qn{6} How many of the following reactions can give Benzene as final product.\par
\qfig{p2_chem_s3_q06_A}
\qn{7}\par
\qfig{p2_chem_s3_q06_8}
Find out degree of unsaturation of product Q :
\qn{8} How many stereoisomers are possible in \qfigi{p2_chem_s3_q08_1}
\partsec{3}{MATHEMATICS}
\sect{SECTION-I(i)}
\qn{1} If x$_1$ and x$_2$ are the roots of the equation $\log_\pi\left((\sin^{-1}(x))(\cos^{-1}(x))\right)=\log_{\sqrt\pi}\left(\dfrac\pi{\sqrt{18}}\right)$ for x $\in$ (0, 1) then the value of $\tan^{-1}\left(x_1^{\,2}+x_2^{\,2}\right)$ is equal to -
\begin{qoptions}
\opt{A}{$\dfrac\pi3$}
\opt{B}{$\dfrac\pi4$}
\opt{C}{$\dfrac\pi6$}
\opt{D}{$\dfrac\pi8$}
\end{qoptions}
\qn{2} Let a$_\mathrm n$ and b$_\mathrm n$ be the sequences of real numbers such that for all integers n $\geq$ 0 where i = $\sqrt{-1}$ and (3 + i)$^\mathrm n$ = (a$_\mathrm n$ + b$_\mathrm n$ i) then the value of $\displaystyle\sum_{n=0}^{\infty}\dfrac{a_nb_n}{12^n}$ is equal to
\begin{qoptions}
\opt{A}{$\dfrac9{13}$}
\opt{B}{$\dfrac5{13}$}
\opt{C}{$\dfrac{12}{13}$}
\opt{D}{$\dfrac7{13}$}
\end{qoptions}
\qn{3} $\displaystyle\int_0^{\pi/3}\dfrac{\{\tan x\}}{\tan x}dx$ is equal to (\{x\} denotes fractional part of x)
\begin{qoptions}
\opt{A}{$\dfrac\pi3+\dfrac12\ell\mathrm n\left(\dfrac32\right)$}
\opt{B}{$\dfrac\pi3-\dfrac12\ell\mathrm n\left(\dfrac34\right)$}
\opt{C}{$\dfrac\pi3-\dfrac12\ell\mathrm n\left(\dfrac32\right)$}
\opt{D}{$\dfrac\pi3+\dfrac12\ell\mathrm n\left(\dfrac34\right)$}
\end{qoptions}
\qn{4} In a cube ABCDEFGH with opposite vertices C and E, J and I are the mid points of segments FB and HD respectively. Let R be the ratio of the area of the cross section EJCI to the area of one of the faces of the cube then value of R$^2$ is equal to -
\begin{qoptions}
\opt{A}{$\dfrac54$}
\opt{B}{$\dfrac32$}
\opt{C}{$\dfrac43$}
\opt{D}{$\dfrac{25}{16}$}
\end{qoptions}
\sect{SECTION-I(ii)}
\qn{5} If $^n$C$_2$ = f(n) and S = $^{f(3)}$C$_2$ + $^{f(4)}$C$_2$ + $^{f(5)}$C$_2$ + ... $^{f(40)}$C$_2$ then which options are correct for S
\begin{qoptions}
\opt{A}{Total number of divisors are 144}
\opt{B}{Total number of odd divisors are 48}
\opt{C}{S is divisible by 78}
\opt{D}{Total number of prime factors are 7}
\end{qoptions}
\qn{6} Let z be a complex number satisfying $12|z|^2=2|z+2|^2+|z^2+1|^2+31$ then which options are correct
\begin{qoptions}
\opt{A}{$\left|z+\dfrac6z\right|=2$}
\opt{B}{Re(z) = $-$1}
\opt{C}{$|z|=\sqrt6$}
\opt{D}{$\dfrac{3\pi}4<\mathrm{Arg}(z)<\pi$}
\end{qoptions}
\qn{7} For regular tetrahedron ABCD, points E,F are on AB, AC respectively, such that BE= 3, EF = 4, and EF is parallel to plane BCD. Then which options are correct
\begin{qoptions}
\opt{A}{DE = $\sqrt{37}$}
\opt{B}{Area of $\Delta$DEF = $2\sqrt{33}$}
\opt{C}{Height of tetrahedron $\sqrt{33}$}
\opt{D}{Volume of tetrahedron = $\dfrac{49}4\sqrt{\dfrac{37}3}$}
\end{qoptions}
\qn{8} For triangle ABC, A(0, 8), B(6, 0) and C(x, y) are three points such that C(x, y) lie on circle (x $-$ 8)$^2$ + (y $-$ 6)$^2$ = 9. If $\Delta$ is the area of triangle ABC and z$_\mathrm m$ \& z$_\mathrm M$ is a point corresponding to complex number of point C for which area is minimum and maximum respectively then select correct options
\begin{qoptions}
\opt{A}{$\Delta_{\max}$ = 41}
\opt{B}{$\Delta_{\min}$ = 11}
\opt{C}{$\arg\left(\dfrac{5\left(z_M+\bar z_M\right)-100}{5\left(z_M-\bar z_M\right)-78}\right)=-\dfrac\pi2$}
\opt{D}{$\arg(z_m)<\dfrac\pi4$}
\end{qoptions}
\qn{9} Let f : N $\longrightarrow$ Q be a function where N denotes the set of natural numbers and Q denotes the set of rational numbers. f(1) = $\dfrac32$ and $f(x+y)=\left(1+\dfrac y{x+1}\right)f(x)+\left(1+\dfrac x{y+1}\right)f(y)+x^2y+xy+xy^2$ for all natural numbers x, y the select correct options :-
\begin{qoptions}
\opt{A}{f(20) = 4305}
\opt{B}{If f(n) = 1 then only one root lie only between (0, 1)}
\opt{C}{3$^2$ + 6$^2$ + 9$^2$ + ... + (3n)$^2$ = 6 f(n)}
\opt{D}{If 2f(n) = (2023) ($^{n+1}$C$_2$) then n = 1011}
\end{qoptions}
\qn{10} Let y(x) be a solution of the differential equation $\left(x-\dfrac{dy}{dx}\right)\sin x=y\cos x$, $x\in\left(0,\dfrac\pi2\right]$. If $y\left(\dfrac\pi2\right)=1$ then choose the correct statement(s) :
\begin{qoptions}
\opt{A}{y(x) attains a maximum value in $\left(0,\dfrac\pi2\right]$}
\opt{B}{y(x) attains a minimum value in $\left(0,\dfrac\pi2\right]$}
\opt{C}{$\displaystyle\lim_{x\to0^+}y(x)=0$}
\opt{D}{$0<y\left(\dfrac\pi4\right)<1$}
\end{qoptions}
\sect{SECTION-III}
\qn{1} For parabola y$^2$ = 4x, let P(11, b) be any point inside the parabola and F be the focus of the parabola. If the point Q on the parabola such that FQ + QP is minimum and R is a point of intersection of axis and directrix. Then the area of parallelogram FPQR is equal to
\qn{2} Let A and B are 2 $\times$ 2 matrice such that tr(A) = 7, det(A) = 22, tr(B) = 4, det(B) = 11, tr(AB) = 33. If det(A + xB) = 272 then the value of x is (x $\in$ I$^+$)
\qn{3} Let n be a positive integer such that exactly one root of the equation $4x^2-\left(4\sqrt3+4\right)x+\sqrt3n-24=0$ is an integer. Then the value of $\dfrac n2$ is
\qn{4} If I = $\displaystyle\int_0^\infty\dfrac1{\left(1+x^{2024}\right)\left(1+x^2\right)}dx$ then the value [I] ([ ] Greatest Integer Function)
\qn{5} Let f : \{1, 2, 3, 4\} $\to$ \{a, b, c, d, e, f\}. If f(1) $\neq$ a or f(2) $\neq$ b or both then number of functions are k then sum of digits of k is
\qn{6} A purchaser of electrical components buys them in lots of size 10. It is his policy to inspect 3 compononts randomly from a lot and to accept the lot only if all 3 are non defective. If 30\% of the lots have x defective components and 70\% have only 1. If purchases reject 46\% of lots then value of x is
\qn{7} Let T$_\mathrm n$ consists of general term of positive numbers (where n is positive integer). If $\dfrac{T_n+2}2=\sqrt{2S_n}$ and S$_\mathrm n$ is sum of n terms, then the value of (T$_{2024}$ $-$ 8090) is equal to
\qn{8} Let f(x) = e$^{-|x|}$, x $\in$ R and g(t) = $\displaystyle\int_{-1}^{1}f\left(\dfrac xt\right)dx$, t $\neq$ 0 then value of $\displaystyle\lim_{t\to0^+}\dfrac{g(t)}t$ is equal to
\end{multicols}
\end{document}
